\section{Geometric applications of integrals}

Integrals can measure geometric quantities by accumulating small pieces. Area, volume, and length can often be expressed as integrals when the object is described by a function.

The main idea is to cut a complicated object into many thin pieces, approximate each piece by something simple, and then pass to a limit. The integral represents this limiting accumulation.

\subsection*{Area under a graph}

If \(f(x)\ge0\) on \([a,b]\), then
\[
\int_a^b f(x)\,dx
\]
gives the area under the graph of \(f\) above the \(x\)-axis.

\begin{examplebox}
A tunnel entrance has the shape of the parabola
\[
y=4-x^2.
\]
Find the area of the entrance above the \(x\)-axis.

The parabola meets the \(x\)-axis when
\[
4-x^2=0,
\]
so
\[
x=-2
\qquad\text{or}\qquad
x=2.
\]
The area is
\[
A=\int_{-2}^{2}(4-x^2)\,dx.
\]
Compute:
\[
A=\left[4x-\frac{x^3}{3}\right]_{-2}^{2}.
\]
At \(x=2\),
\[
4x-\frac{x^3}{3}=8-\frac83=\frac{16}{3}.
\]
At \(x=-2\),
\[
4x-\frac{x^3}{3}=-8+\frac83=-\frac{16}{3}.
\]
Thus
\[
A=\frac{16}{3}-\left(-\frac{16}{3}\right)=\frac{32}{3}.
\]
\end{examplebox}

The limits of integration are not arbitrary. They come from where the curve meets the axis.

\subsection*{Area of a circle}

The upper semicircle of radius \(R\) centered at the origin is described by
\[
y=\sqrt{R^2-x^2},
\qquad -R\le x\le R.
\]
The area of the full circle is twice the area under this upper semicircle:
\[
A=2\int_{-R}^{R}\sqrt{R^2-x^2}\,dx.
\]
This integral evaluates to
\[
\pi R^2.
\]
Thus the familiar area formula for a circle can be understood as accumulated vertical slices.

\begin{remarkbox}
Even familiar formulas can be reinterpreted by integration. The formula \(A=\pi R^2\) is not only memorized geometry; it can also be derived from accumulation.
\end{remarkbox}

\subsection*{Volumes of solids of revolution}

If a non-negative function \(f\) is rotated around the \(x\)-axis, cross-sections perpendicular to the \(x\)-axis are disks of radius \(f(x)\). The area of each disk is
\[
\pi[f(x)]^2.
\]
Thus the volume is
\[
V=\int_a^b \pi[f(x)]^2\,dx.
\]

\begin{examplebox}
A vessel is formed by rotating
\[
y=\sqrt{x}
\]
around the \(x\)-axis for \(0\le x\le4\). Its volume is
\[
V=\int_0^4 \pi(\sqrt{x})^2\,dx.
\]
Since \((\sqrt{x})^2=x\),
\[
V=\pi\int_0^4 x\,dx
=\pi\left[\frac{x^2}{2}\right]_0^4
=8\pi.
\]
\end{examplebox}

The formula is simple because each thin slice is approximately a disk.

\subsection*{Arc length}

The length of a curve \(y=f(x)\) on \([a,b]\) is given by
\[
L=\int_a^b\sqrt{1+[f'(x)]^2}\,dx.
\]
This formula combines derivatives and integrals. The derivative measures local slope; the integral accumulates small pieces of length.

\begin{examplebox}
Suppose a hanging rope is approximated by
\[
f(x)=\frac{1}{10}x^2,
\qquad -10\le x\le10.
\]
Then
\[
f'(x)=\frac{x}{5}.
\]
The rope length is
\[
L=\int_{-10}^{10}\sqrt{1+\frac{x^2}{25}}\,dx.
\]
This integral may require a more advanced method to evaluate exactly, but the important modelling step is already clear: slope determines local stretching, and the integral accumulates length.
\end{examplebox}

\begin{summarybox}
For geometric integral problems, ask:
\begin{itemize}
\item What small pieces are being accumulated?
\item What are the correct bounds of integration?
\item Is the integral measuring area, volume, or length?
\item Is the quantity signed or ordinary geometric size?
\item Does the formula come from vertical slices, disks, or arc-length pieces?
\end{itemize}
\end{summarybox}

Geometric applications show that integration is a method for building whole quantities from infinitesimal pieces.